\subsection{Server-Side Attacks}

\begin{itemize}
\item \textbf{Binary Planting:} Uploading and executing unauthorized binary files on the server, usually a specific location. Binary planting can be used to escalate privileges. An example of this is the Stuxnet virus using the  Microsoft Windows Print Spooler Service Remote Code Execution Vulnerability (CVE-2010-2729 / BID 43073). Also known as \textbf{DLL Planting} or \textbf{DLL Hijacking}.
\item \textbf{SQL Injection}
\begin{itemize}
    \item Blind SQL Injection: Extracting data from a database by sending queries and observing the server's response time or behavior.
\end{itemize}
\item Blind XPath Injection: Extracting data from an XML document by sending queries and observing the server's response time or behavior.
\item Buffer Overflow via Environment Variables: Manipulating environment variables to cause a buffer overflow and execute arbitrary code.
\item Buffer Overflow Attack: Overwriting memory beyond the bounds of a buffer to corrupt data or execute arbitrary code.
\item \textbf{CORS}
\begin{itemize}
    \item CORS Origin Header Scrutiny: Exploiting misconfigured Cross-Origin Resource Sharing (CORS) headers to access restricted resources.
    \item CORS Request Preflight Scrutiny: Exploiting misconfigured CORS preflight requests to bypass security checks.
\end{itemize}
\item \textbf{File Injection:}
\begin{itemize}
    \item CSS 
    \item CSV Injection: Injecting malicious payloads into CSV files to execute arbitrary code or perform unauthorized actions.
\end{itemize}

\item Cache Poisoning: Manipulating cached data to serve malicious content to users.
\item Cash Overflow: Manipulating financial calculations to cause integer overflows and perform unauthorized transactions.
\item Code Injection: Injecting malicious code into a vulnerable application to execute arbitrary commands or manipulate data.
\item Command Injection: Executing arbitrary commands on the server by injecting them into vulnerable input fields.
\item Comment Injection Attack: Injecting malicious comments or code into an application's data store or output.
\item Content Spoofing: Tricking users into believing that malicious content is legitimate by mimicking the appearance of a trusted website.
\item Custom Special Character Injection: Injecting special characters into input fields to bypass validation or cause unintended behavior.
\item Denial of Service (DoS): Overwhelming a server with requests to exhaust its resources and make it unavailable to users.
\item Direct Dynamic Code Evaluation - Eval Injection: Injecting malicious code into an application's eval() function to execute arbitrary commands.
\item Embedding Null Code: Injecting null bytes into input fields to bypass validation or manipulate file paths.
\item Execution After Redirect (EAR): Executing malicious code after a redirect to bypass security checks.
\item Forced Browsing: Accessing restricted pages or resources by guessing or manipulating URLs.
\item Form Action Hijacking: Modifying the action attribute of HTML forms to submit data to an attacker-controlled server.
\item Format String Attack: Exploiting format string vulnerabilities to read or write arbitrary memory locations.
\item Full Path Disclosure: Revealing the full path of files or directories on the server.
\item Function Injection: Injecting malicious code into a vulnerable function to execute arbitrary commands or manipulate data.
\item HTTP Response Splitting: Injecting newline characters into HTTP headers to split the response and insert malicious content.
\item LDAP Injection: Injecting malicious LDAP queries to extract sensitive information or manipulate directory services.
\item Log Injection: Injecting malicious payloads into application logs to execute arbitrary code or manipulate log data.
\item Server-Side Includes (SSI) Injection: Injecting SSI directives to execute arbitrary commands or include malicious files.
\item Server-Side Request Forgery (SSRF): Tricking a server into making unauthorized requests to internal or external resources.
\item Session Prediction: Guessing or predicting valid session identifiers to hijack user sessions.
\item Session Fixation: Forcing a user to use a predefined session identifier to hijack their session.
\item Session Hijacking Attack: Stealing or predicting session identifiers to take over user sessions.
\item Setting Manipulation: Modifying application settings or configuration files to alter its behavior or gain unauthorized access.
\item Special Element Injection: Injecting special elements or tags into input fields to bypass validation or manipulate the application's behavior.
\item Unicode Encoding: Exploiting Unicode encoding vulnerabilities to bypass input validation or inject malicious payloads.
\item Web Parameter Tampering: Manipulating URL parameters or form data to perform unauthorized actions or bypass security checks.
\item Windows ::DATA Alternate Data Stream: Hiding malicious files or data in alternate data streams on Windows systems.
\item XML External Entity (XXE) Injection: Exploiting XML parsers to access unauthorized files or network resources.
\item XPath Injection: Injecting malicious XPath queries to extract sensitive information from XML documents.
\item Cross-Site Request Forgery (CSRF): Tricking users into performing unintended actions on a web application in which they are authenticated.
\item Web Service Amplification Attack: Exploiting web services to amplify denial of service attacks by sending large amounts of data.
\end{itemize}


\hypertarget{other-attacks}{%
\subsubsection{Other Attacks:}\label{other-attacks}}

\begin{itemize}
\item
  Brute Force Attack
\item
  Cross-User Defacement
\item
  Forced browsing
\item
  Form action hijacking by Robert Gilbert (amroot)
\item
  Man-in-the-browser attack
\item
  Manipulator-in-the-middle attack
\item
  Parameter Delimiter
\item
  Path Traversal
\item
  Regular expression Denial of Service - ReDoS by Adar Weidman
\item
  Repudiation Attack
\item
  Resource Injection
\item
\begin{verbatim}
* **Inband**. The same channel is used to inject SQL code and extract sensitive data.
\end{verbatim}
\item
\begin{verbatim}
* **Out-of-band.**  Data is retrieved using a different channel.
\end{verbatim}
\item
\begin{verbatim}
* **Inferential.**  Gaining information from behaviour server based on sent requests.
\end{verbatim}
\item
\begin{verbatim}
* Tools:
\end{verbatim}
\item
\begin{verbatim}
    * mieliekoek.pl
\end{verbatim}
\item
\begin{verbatim}
    * wpoison
\end{verbatim}
\item
\begin{verbatim}
    * sqlmap
\end{verbatim}
\item
\begin{verbatim}
    * wapiti
\end{verbatim}
\item
\begin{verbatim}
    * w3af
\end{verbatim}
\item
\begin{verbatim}
    * paros
\end{verbatim}
\item
\begin{verbatim}
    * sqid
\end{verbatim}
\item
\begin{verbatim}
* User input
\end{verbatim}
\item
\begin{verbatim}
* Cookies
\end{verbatim}
\item
\begin{verbatim}
* HTTP Headers
\end{verbatim}
\item
\begin{verbatim}
* Second-order SQL Injections
\end{verbatim}
\item
  Traffic flood
\item
  Trojan Horse
\item
  Unicode Encoding
\item
  Windows ::DATA Alternate Data Stream
\end{itemize}